\textbf{Optimiser budget of the baselines.} Because the ESN-versus-MLP comparison could depend on how long the MLP is trained, Table~\ref{tab:steps} and Fig.~\ref{fig:steps} vary the number of Adam steps of the MLP and GRU on the test seeds 0--2, with everything else at the tuned values (this sweep has no hypothesis and was not used for tuning). On Lorenz-63 the MLP's mean rises from 4000 to 16000 steps (\rhval{sweep_steps/mlp-delay@steps=4000/lorenz63/vpt/mean:2} to \rhval{sweep_ntrain/mlp-delay@n_train=5000/lorenz63/vpt/mean:2}) and little beyond (\rhval{sweep_steps/mlp-delay@steps=32000/lorenz63/vpt/mean:2} at 32000); the ESN/MLP ratio falls from \rhval{cmp/st_mlp4000/mlp-delay/lorenz63/vpt/ratio:2} to \rhval{cmp/st_mlp32000/mlp-delay/lorenz63/vpt/ratio:2}, so the ESN stays ahead at every budget tested but the size of the gap depends on the budget. On Lorenz-96 the MLP's mean rises at every doubling, to \rhval{sweep_steps/mlp-delay@steps=32000/lorenz96/vpt/mean:2} at 32000 steps; with 3 seeds we do not claim that the individual steps beyond the tuned 8000 are significant, but the trend suggests that Table~\ref{tab:main} understates, not overstates, the MLP's lead over the ESN on this task. The GRU's mean also rises with the budget on both tasks; we make no significance claim for its differences to the tuned value, and the ESN/GRU ratio stays at or above \rhval{cmp/st_gru12000/gru/lorenz96/vpt/ratio:2}.

\begin{table*}[t]\centering\caption{VPT versus number of optimiser steps (seeds 0--2). $\ast$: the tuned budget (the main runs of these seeds); ESN/sys.: ratio of the ESN's mean VPT (same seeds, Table~\ref{tab:ntrain}, row 5000) to this row's (\texttt{rh compare}).}\label{tab:steps}
\small\begin{tabular}{lrcccccc}\toprule
System & steps & L63 VPT & L63 ESN/sys. & L63 fit (s) & L96 VPT & L96 ESN/sys. & L96 fit (s)\\
\midrule
MLP & 4000 & 3.51 $\pm$ 0.65 & \rhval{cmp/st_mlp4000/mlp-delay/lorenz63/vpt/ratio:2} & 6.0 & 2.38 $\pm$ 0.12 & \rhval{cmp/st_mlp4000/mlp-delay/lorenz96/vpt/ratio:2} & 6.5\\
MLP & 8000 & 5.03 $\pm$ 0.94 & \rhval{cmp/st_mlp8000/mlp-delay/lorenz63/vpt/ratio:2} & 14.1 & 2.72 $\pm$ 0.09$^{\ast}$ & \rhval{cmp/st_mlp8000/mlp-delay/lorenz96/vpt/ratio:2} & 15.9\\
MLP & 16000 & 5.55 $\pm$ 0.24$^{\ast}$ & \rhval{cmp/st_mlp16000/mlp-delay/lorenz63/vpt/ratio:2} & 28.5 & 3.13 $\pm$ 0.18 & \rhval{cmp/st_mlp16000/mlp-delay/lorenz96/vpt/ratio:2} & 20.3\\
MLP & 32000 & 5.68 $\pm$ 0.17 & \rhval{cmp/st_mlp32000/mlp-delay/lorenz63/vpt/ratio:2} & 55.7 & 3.49 $\pm$ 0.37 & \rhval{cmp/st_mlp32000/mlp-delay/lorenz96/vpt/ratio:2} & 53.5\\
\midrule
GRU & 3000 & 2.88 $\pm$ 0.23 & \rhval{cmp/st_gru3000/gru/lorenz63/vpt/ratio:2} & 20.9 & 1.03 $\pm$ 0.10 & \rhval{cmp/st_gru3000/gru/lorenz96/vpt/ratio:2} & 21.6\\
GRU & 6000 & 3.60 $\pm$ 0.72$^{\ast}$ & \rhval{cmp/st_gru6000/gru/lorenz63/vpt/ratio:2} & 56.5 & 1.20 $\pm$ 0.20$^{\ast}$ & \rhval{cmp/st_gru6000/gru/lorenz96/vpt/ratio:2} & 62.5\\
GRU & 12000 & 3.80 $\pm$ 0.56 & \rhval{cmp/st_gru12000/gru/lorenz63/vpt/ratio:2} & 107.4 & 1.31 $\pm$ 0.26 & \rhval{cmp/st_gru12000/gru/lorenz96/vpt/ratio:2} & 112.4\\
\bottomrule\end{tabular}
\end{table*}
\begin{figure}[t]\centering\includegraphics[width=\columnwidth]{sweep_steps_fig.pdf}
\caption{VPT of the MLP and GRU versus optimiser steps (log axis), seeds 0--2.}\label{fig:steps}\end{figure}

\textbf{Component ablation.} Table~\ref{tab:abl} switches off the ESN's squared readout features and the GRU's input noise (5 seeds, the data and settings of Table~\ref{tab:main}, other hyperparameters unchanged); the tests are in Table~\ref{tab:cmp}. The tuned MLP uses no input noise, so it has no ablation. Without squared features the ESN loses VPT on both tasks (\rhval{main/esn-without-squared-features/lorenz63/vpt/mean:2} versus \rhval{main/esn/lorenz63/vpt/mean:2} on Lorenz-63, $p=\rhval{cmp/main/esn-without-squared-features/lorenz63/vpt/paired_p:3}$; \rhval{main/esn-without-squared-features/lorenz96/vpt/mean:2} versus \rhval{main/esn/lorenz96/vpt/mean:2} on Lorenz-96, $p=\rhval{cmp/main/esn-without-squared-features/lorenz96/vpt/paired_p:3}$). On Lorenz-63 its \texttt{clim\_ok} is unchanged (\rhval{main/esn-without-squared-features/lorenz63/clim_ok/mean:2} versus \rhval{main/esn/lorenz63/clim_ok/mean:2}), but on Lorenz-96 every free run leaves the attractor region (\texttt{blowup}$=1$). As defined in Sec.~\ref{sec:method}, such a run exceeds three times the reference range and need not be numerically unbounded. The GRU without input noise keeps its VPT on both tasks ($p=\rhval{cmp/abl_gru/gru-without-input-noise/lorenz63/vpt/paired_p:3}$ and $\rhval{cmp/abl_gru/gru-without-input-noise/lorenz96/vpt/paired_p:3}$) but on Lorenz-96 its free runs leave the attractor region in \rhval{main/gru-without-input-noise/lorenz96/blowup/mean:2} of the cases (\texttt{clim\_ok} \rhval{main/gru-without-input-noise/lorenz96/clim_ok/mean:2} versus \rhval{main/gru/lorenz96/clim_ok/mean:2}). On Lorenz-63 its \texttt{clim\_ok} is lower in mean (\rhval{main/gru-without-input-noise/lorenz63/clim_ok/mean:2} versus \rhval{main/gru/lorenz63/clim_ok/mean:2}) but the difference is not significant ($p=\rhval{cmp/abl_gru/gru-without-input-noise/lorenz63/clim_ok/paired_p:3}$). So on Lorenz-96 input noise matters for the GRU's long rollouts and not for its short-term accuracy: a case where VPT and climate fidelity separate within one model family.

\begin{table}[t]\centering\caption{Component ablation (5 seeds). \texttt{blowup}: fraction of the 20 free runs per seed that leave the attractor region.}\label{tab:abl}
\small\setlength{\tabcolsep}{4pt}\begin{tabular}{llcccc}
\toprule
Method & Task & vpt $\uparrow$ & clim\_ok $\uparrow$ & blowup $\downarrow$ & $n$ \\
\midrule
ESN & L63 & 9.55 $\pm$ 0.32 & 0.85 $\pm$ 0.14 & 0.00 $\pm$ 0.00 & 5 \\
ESN no squares & L63 & 8.13 $\pm$ 0.75 & 0.84 $\pm$ 0.12 & 0.00 $\pm$ 0.00 & 5 \\
GRU & L63 & 3.79 $\pm$ 0.61 & 0.86 $\pm$ 0.11 & 0.00 $\pm$ 0.00 & 5 \\
GRU no noise & L63 & 3.67 $\pm$ 0.55 & 0.63 $\pm$ 0.35 & 0.00 $\pm$ 0.00 & 5 \\
\midrule
ESN & L96 & 2.11 $\pm$ 0.34 & 1.00 $\pm$ 0.00 & 0.00 $\pm$ 0.00 & 5 \\
ESN no squares & L96 & 0.96 $\pm$ 0.11 & 0.00 $\pm$ 0.00 & 1.00 $\pm$ 0.00 & 5 \\
GRU & L96 & 1.33 $\pm$ 0.23 & 1.00 $\pm$ 0.00 & 0.00 $\pm$ 0.00 & 5 \\
GRU no noise & L96 & 1.33 $\pm$ 0.20 & 0.01 $\pm$ 0.02 & 0.99 $\pm$ 0.02 & 5 \\
\bottomrule
\end{tabular}
\end{table}

\textbf{Spectral radius (H3).} Table~\ref{tab:rho} sweeps $\rho$ with all else at the tuned values (seeds 0--2; the 0.4 row is the main runs of those seeds). The best seed-mean VPT is \rhval{sweep_rho/esn@rho=0.1/lorenz63/vpt/mean:2} on Lorenz-63 and \rhval{sweep_rho/esn@rho=0.1/lorenz96/vpt/mean:2} on Lorenz-96 (both at $\rho=0.1$) and the worst \rhval{sweep_rho/esn@rho=2.0/lorenz63/vpt/mean:2} and \rhval{sweep_rho/esn@rho=2.0/lorenz96/vpt/mean:2} (both at $\rho=2$); on each task the best is more than twice the worst, so H3 is supported. The shapes differ. On Lorenz-63 VPT changes little for $\rho\le0.7$ and falls for $\rho\ge1.0$, while \texttt{clim\_ok} shows no trend (\rhval{sweep_rho/esn@rho=0.7/lorenz63/clim_ok/mean:2} to \rhval{sweep_rho/esn@rho=1.4/lorenz63/clim_ok/mean:2} at every radius, with seed stds of \rhval{sweep_rho/esn@rho=none/lorenz63/clim_ok/std:2} to \rhval{sweep_rho/esn@rho=2.0/lorenz63/clim_ok/std:2}): there the short-term metric degrades to less than half of its best value without a visible change in the climate metric. On Lorenz-96 the mean VPT falls monotonically from the smallest radius tested and \texttt{clim\_ok} is 0 for $\rho\ge1.0$. The difference between $\rho=0.1$ and the tuned 0.4 is within the seed std of the 0.4 row on both tasks, so the sweep does not show that the main ESN row is sub-optimal; it does show that the tuning grid, which started at 0.4, did not bracket the optimum. Good performance at small radii agrees with reports that reservoirs with a small spectral radius reconstruct attractors better~\cite{hart2023attractor} and that reservoirs with zero spectral radius can forecast Lorenz-63~\cite{griffith2019forecasting}.

\begin{table}[t]\centering\caption{ESN VPT and \texttt{clim\_ok} versus spectral radius $\rho$ (top; the 0.4 row is the tuned setting) and versus reservoir size $N$ (bottom). Mean $\pm$ std, seeds 0--2.}\label{tab:rho}
\small\setlength{\tabcolsep}{3pt}\begin{tabular}{rcccc}\toprule
$\rho$ & L63 VPT & L96 VPT & L63 clim\_ok & L96 clim\_ok\\
\midrule
0.1 & 9.70 $\pm$ 0.58 & 2.30 $\pm$ 0.15 & 0.83 $\pm$ 0.16 & 1.00 $\pm$ 0.00\\
0.4 & 9.60 $\pm$ 0.31 & 2.04 $\pm$ 0.46 & 0.85 $\pm$ 0.13 & 1.00 $\pm$ 0.00\\
0.7 & 9.42 $\pm$ 0.30 & 1.38 $\pm$ 0.13 & 0.78 $\pm$ 0.16 & 0.88 $\pm$ 0.08\\
1.0 & 8.30 $\pm$ 0.41 & 0.64 $\pm$ 0.13 & 0.80 $\pm$ 0.20 & 0.00 $\pm$ 0.00\\
1.4 & 4.06 $\pm$ 0.49 & 0.28 $\pm$ 0.07 & 0.87 $\pm$ 0.19 & 0.00 $\pm$ 0.00\\
2.0 & 2.32 $\pm$ 0.84 & 0.21 $\pm$ 0.04 & 0.78 $\pm$ 0.25 & 0.00 $\pm$ 0.00\\
\bottomrule\end{tabular}


\vspace{3pt}\begin{tabular}{rcccc}\toprule
$N$ & L63 VPT & L96 VPT & L63 clim\_ok & L96 clim\_ok\\
\midrule
100 & 6.80 $\pm$ 0.45 & 0.34 $\pm$ 0.04 & 0.90 $\pm$ 0.09 & 0.00 $\pm$ 0.00\\
250 & 8.78 $\pm$ 0.13 & 1.42 $\pm$ 0.21 & 0.85 $\pm$ 0.09 & 1.00 $\pm$ 0.00\\
500 & 9.60 $\pm$ 0.31 & 2.04 $\pm$ 0.46 & 0.85 $\pm$ 0.13 & 1.00 $\pm$ 0.00\\
1000 & 9.44 $\pm$ 0.13 & 2.62 $\pm$ 0.30 & 0.87 $\pm$ 0.15 & 1.00 $\pm$ 0.00\\
2000 & 9.93 $\pm$ 0.03 & 3.28 $\pm$ 0.57 & 0.87 $\pm$ 0.08 & 1.00 $\pm$ 0.00\\
\bottomrule\end{tabular}
\end{table}

\textbf{Reservoir size (H4).} In Table~\ref{tab:rho} (bottom) the mean VPT at $N=1000$ exceeds that at $N=100$ on both tasks (\rhval{sweep_size/esn@size=1000/lorenz63/vpt/mean:2} versus \rhval{sweep_size/esn@size=100/lorenz63/vpt/mean:2} on Lorenz-63; \rhval{sweep_size/esn@size=1000/lorenz96/vpt/mean:2} versus \rhval{sweep_size/esn@size=100/lorenz96/vpt/mean:2} on Lorenz-96), so H4 is supported. On Lorenz-63 the gain is concentrated below $N=500$; on Lorenz-96 the mean rises at every size up to \rhval{sweep_size/esn@size=2000/lorenz96/vpt/mean:2} at $N=2000$. At $N=100$ on Lorenz-96 \texttt{clim\_ok} is 0, and it is 1 at all larger sizes; on Lorenz-63 it lies between \rhval{sweep_size/esn@size=250/lorenz63/clim_ok/mean:2} and \rhval{sweep_size/esn@size=100/lorenz63/clim_ok/mean:2} at every size.
